\documentclass[11pt,a4paper]{article}
\usepackage[T1]{fontenc}\usepackage{amsmath,amssymb}\usepackage[margin=2.5cm]{geometry}
\title{Phase Ordering Described by the Allen--Cahn Equation}\author{I. Phase and J. Field}\date{}
\begin{document}\maketitle
\section*{1. Model}
The evolution of a non-conserved order parameter $c(x,y,t)$ in the unit square $\Omega=(0,1)^2$ is modelled as an $L^2$ gradient flow of the Ginzburg--Landau free energy

\begin{equation}
F[c] = \int_\Omega \left( f(c) + \frac{\kappa}{2}\,|\nabla c|^2 \right) d\Omega ,
\end{equation}
with the double-well potential

\begin{equation}
f(c) = \frac14\left(c^2-1\right)^2 .
\end{equation}
The resulting Allen--Cahn equation reads

\begin{equation}
\frac{\partial c}{\partial t} = -M\left( f'(c) - \kappa\,\nabla^2 c \right) \quad \text{in } \Omega,\ t>0 ,
\end{equation}
where $M$ is the mobility and $\kappa$ the gradient-energy coefficient.

\section*{2. Boundary and initial conditions}
No flux crosses the boundary of the domain,

\begin{equation}
\nabla c\cdot n = 0 \quad \text{on } \partial\Omega ,
\end{equation}
and the initial state is a smooth saddle-shaped field between the two wells,

\begin{equation}
c(x,y,0) = \tfrac12\,(2x-1)(2y-1) .
\end{equation}
\section*{3. Parameters}
We use $M=1$ and $\kappa=0.01$. The simulation runs up to $t=1$ with the time step $\Delta t=0.01$.

\section*{4. Discussion}
Because the flow is not conservative the mean of $c$ is not preserved, whereas the free energy $F$ decreases monotonically in time.

\end{document}